% Arabic Abstract (XeLaTeX/polyglossia)
\begin{Arabic}
\chapter*{المستخلص}

يربط التوليد المعزَّز بالاسترجاع (RAG) مخرجات النماذج اللغوية الكبيرة (LLM) بمعرفة خارجية، غير أن استرجاع المعلومات العربية يُخفق في كثير من الاستعلامات بسبب الثراء الصرفي والتباين الإملائي وعدم التطابق المعجمي بين الاستعلام والمستند. ويُقدِّم تحسين الاستعلام (QE)، أي تعديل الاستعلام قبل الاسترجاع، حلاً يمكن إضافته دون تغيير بقية النظام، إلا أن تقنياته وُضعت أساساً للغة الإنجليزية معتمدةً على نماذج احتكارية أكبر بكثير، ولم تثبت بعدُ فاعلية تطبيقها على العربية ولا موضعها الأمثل داخل بنية هجينة. تبحث هذه الأطروحة مدى قدرة تحسين الاستعلام المعتمد على النماذج اللغوية الكبيرة، بنوعيه الأعمى والموجَّه بالمدوَّنة المستهدفة، على الارتقاء باسترجاع المعلومات العربية في أنماط الاسترجاع الثلاثة: المتناثر والكثيف والهجين، مع الاقتصار على نماذج متاحة للعموم. وتمثلت أهدافها في وضع خطوط أساس، وتكييف التحسين التوليدي للعربية، والمقارنة بين النماذج التوليدية، وتحديد موضع التوسيع داخل النظام الهجين. أُجريت التجارب على الجزء العربي من معيار MIRACL، الذي يضم 2,896 استعلاماً و2.06 مليون مقطع نصي، باستخدام المسترجع الكثيف mDPR والمسترجع المتناثر BM25S. وأظهر تحليل الأخطاء إخفاق 34\% من الاستعلامات، فدفع ذلك إلى تكييف تقنية Query2Doc للعربية. وقورنت عشرة نماذج لغوية توليدية، وأثبتت التجارب صلاحية تسعة منها، ثم اختُبرت طرائق تكرار الاستعلام والدمج الهجين والتوسيع الموجَّه بالمدوَّنة (CSQE)، الذي يسترجع أولاً عدداً قليلاً من المستندات الأعلى ترتيباً من المدوَّنة نفسها، ثم يبني التوسيع من ألفاظها لا من المعرفة الكامنة في النموذج. وقد حسَّن توسيع الاستعلام الاسترجاعَ الكثيف في النماذج التسعة جميعها، غير أنه أضعف الاسترجاع المتناثر في معظمها لأنه يخفف أوزان مصطلحات الاستعلام الأصلية، وهو أثر أزاله تكرار الاستعلام. وتبيَّن أن موضع تطبيق التوسيع عامل حاسم؛ إذ تفوَّق تطبيقه على المسترجع المتناثر وحده على تطبيقه على المسترجعَين معاً أو على المسترجع الكثيف وحده، لأن حجب التوسيع عن المسترجع الكثيف يحفظ التكامل بين المسترجعَين الذي يقوم عليه الدمج. وبذلك رفع النظام النهائي الكسب التراكمي المخصوم المُعيَّر عند الرتبة 10 (NDCG@10) من 0.4621 باستخدام BM25 وحده إلى 0.7137، بزيادة نسبية بلغت 54.5\%، وتجاوز أقوى نظام هجين بلا تحسين للاستعلام بنسبة 13.9\%، باستخدام نموذج توليدي متاح للعموم يضم 8 مليارات معلمة. ومن ثَمَّ فإن تحسين الاستعلام المعتمد على النماذج اللغوية الكبيرة فعّال في أنماط الاسترجاع الثلاثة جميعها، شريطة مواءمة صيغته وموضعه لطبيعة المسترجع.
\end{Arabic}
